%% kapa: exam=1 round=37 year=2026 period=2 form=A subject=appraisal_law subject_ko=감정평가관계법규 seq=3 number=3
%% source: Q-Net 감정평가사 기출문제 (cst00309) artlSeq=5258924 file=raw/1008/5258924/2026년 제37회 감정평가사 1차 2교시.pdf
%% conversion: kordoc 4.24.51 -> tools/gen1.py (auto); review=proofread-sonnet; flags=none
\begin{question}{3}
국토의 계획 및 이용에 관한 법령상 국토교통부장관이 도시ㆍ군관리계획(지구단위계획, 도시혁신계획 및 복합용도계획은 제외)을 중앙도시계획위원회의 심의를 거치지 않고 변경할 수 있는 경우는? (단, 주어진 조건 외에 다른 조건은 고려하지 않음)
\begin{choices}
  \item 도시지역의 확대에 따른 용도지역의 변경인 경우
  \item 단위 도시ㆍ군계획시설부지 면적의 10퍼센트 미만의 변경인 경우
  \item 건축물의 건축에 따른 변속차로의 설치를 위한 도시ㆍ군계획시설의 변경인 경우
  \item 도시지역외의 지역에서 ｢산지관리법｣에 의한 보전산지를 자연환경보전지역으로 결정하는 경우
  \item ｢자연공원법｣에 따른 공원구역을 보전관리지역으로 결정하는 경우
\end{choices}
\end{question}
\answer{3}
